\documentclass[../../../include/open-logic-section]{subfiles}

\begin{document}

\olfileid{sth}{z}{infinity-again}
\olsection{અનંતતા}

અનંત સંખ્યામાં ગણો હોય (જો કોઈ ગણ હોય તો) તેની ખાતરી કરવા માટેની
પૂરતી સ્વયંસિદ્ધિઓ આપણી પાસે પહેલેથી છે. ધારો કે કોઈ ગણ અસ્તિત્વ
ધરાવે છે; ત્યારે $\emptyset$ પણ અસ્તિત્વ ધરાવે છે
(\olref[sth][z][sep]{prop:emptyexists} મુજબ). હવે દરેક ગણ~$x$ માટે
$x \cup \{x\}$ ગણ \olref[sth][z][pairs]{prop:pairsconsequences} મુજબ
અસ્તિત્વ ધરાવે છે. તેથી આ ક્રિયા થોડાં વખત કરવાથી આપણને નીચેના ગણો
મળશે:
\begin{enumerate}
	\item[0.] $\emptyset$ %& $0$ & stage $0$\\
	\item[1.] $ \{\emptyset\}$ %& $1$ & stage $1$\\
	\item[2.] $\{\emptyset, \{\emptyset\}\}$ %& $2$ & stage $2$\\
	\item[3.] $\{\emptyset, \{\emptyset\}, \{\emptyset, \{\emptyset\}\}\}$ %&$3$ & stage $3$\\
	\item[4.] $\{\emptyset, \{\emptyset\}, \{\emptyset, \{\emptyset\}\}, \{\emptyset, \{\emptyset\}, \{\emptyset, \{\emptyset\}\}\}\}$% & $4$ & stage $4$ 
\end{enumerate}
આ બધા ગણો એકબીજાથી જુદા છે તે આપણે તપાસી શકીએ છીએ.

આપણે અમુક કારણોસર ગણતરી $0$થી શરૂ કરી છે. તેમાંનું એક કારણ આ છે.
જે ગણને આપણે “$n$” સંજ્ઞા આપી છે તેના બરાબર $n$ સભ્યો હોય છે,
અને (સહજ વિચાર મુજબ) તે $n$મા તબક્કે રચાય છે, તે તપાસવું બહુ અઘરું
નથી.

પણ અહીં એક ભેદ છે. આનાથી \emph{અનંત સંખ્યામાં} ગણો મળે છે; પરંતુ
કોઈ \emph{અનંત ગણ}, એટલે કે અનંત સંખ્યામાં સભ્યો ધરાવતો ગણ, હોય
તેની ખાતરી થતી નથી. અને આ ભેદ ખૂબ મહત્ત્વનો છે: કોઈ
(ડેડેકિન્ડ) અનંત ગણ ન મળે ત્યાં સુધી આપણે ડેડેકિન્ડ બીજગણિત રચી
શકતા નથી. પણ આપણે ડેડેકિન્ડ બીજગણિત જોઈએ છે, જેથી તેને પ્રાકૃતિક
સંખ્યાઓનો ગણ ગણી શકીએ.
(\olref[sfr][infinite][dedekindsproof]{sec} સાથે સરખાવો.)

અત્યાર સુધી આપેલી સ્વયંસિદ્ધિઓ કોઈ અનંત ગણના અસ્તિત્વની ખાતરી
\emph{આપતી નથી}. તેથી આપણે નવી સ્વયંસિદ્ધિ આપવી પડશે:

\begin{axiom}[અનંતતા]
એવો ગણ $I$ છે કે $\emptyset \in I$, અને જ્યારે પણ $x \in I$ હોય
ત્યારે $x \cup \{x\} \in I$ હોય.
\begin{align*}
	\exists I( & (\exists o \in I)\forall x\ x \notin o \land {}\\
	& (\forall x \in I)(\exists s \in I)\forall z(z \in s \liff (z \in x \lor z = x)))
\end{align*}
\end{axiom}

અનંતતાની સ્વયંસિદ્ધિથી મળતો ગણ $I$ ડેડેકિન્ડ-અનંત છે તે જોઈ શકાય
છે. તેનો વિશિષ્ટ ઘટક $\emptyset$ છે, અને $I$ પરની અનુગામી ક્રિયા
$s(x) = x\cup
\{x\}$ દ્વારા અપાય છે. હવે
\olref[sfr][infinite][dedekind]{thm:DedekindInfiniteAlgebra}માં
ડેડેકિન્ડ-અનંત ગણમાંથી ડેડેકિન્ડ બીજગણિત કેવી રીતે મેળવવું તે
બતાવ્યું હતું; અને તેને આપણે પ્રાકૃતિક સંખ્યાઓનો ગણ ગણીશું.
વધુ ચોક્કસ રીતે:
\begin{defn}\ollabel{defnomega}
અનંતતાની સ્વયંસિદ્ધિથી મળતો કોઈ પણ ગણ $I$ લો. $s$ એ
$s(x) = x \cup \{x\}$ વિધેય છે એમ લો. $\omega =
\closureofunder{s}{\emptyset}$ લો. $\omega$ના સભ્યોને આપણે
\emph{પ્રાકૃતિક સંખ્યાઓ} કહીએ છીએ, અને $n$ એ $\emptyset$ પર
$s$ને $n$ વખત લગાડવાનું પરિણામ છે એમ કહીએ છીએ.
\end{defn}
\sourcecorrection{OLMD-114}{મૂળમાં દરેક આવા $I$ પર $s$ને એક-એક
વિધેય ગણ્યો છે. સભ્યતાનાં ચક્રોને બાકાત રાખતી સ્વયંસિદ્ધિ હજી
અપાઈ નથી, તેથી આ સામાન્ય દાવો અહીં ઉપલબ્ધ સ્વયંસિદ્ધિઓથી મળતો
નથી. અહીં જરૂરી એક-એક વિધેય એ $s$નું $\omega$ પરનું મર્યાદિત
વિધેય છે; $\omega$ સૌથી નાનો અનુગામી-સંવૃત ગણ હોવાથી તેના પર
ગાણિતિક અનુમાનથી અનુગામીઓ ભિન્ન હોવાનું અને $\emptyset$ કોઈનો
અનુગામી ન હોવાનું સ્થાપિત થાય છે. $I$ના બાકીના સભ્યોને તેમના
પોતાના પર મોકલીને આ વિધેયને $I$ પર વિસ્તારી શકાય છે; વિસ્તૃત
વિધેય એક-એક છે અને તેના વિસ્તારમાં $\emptyset$ નથી. તેથી $I$
ડેડેકિન્ડ-અનંત હોવાનો દાવો અને આગળની વ્યાખ્યા યથાવત્ રહે છે.}

હવે થોડાં ફકરા પહેલાં “$n$” સંજ્ઞા આપેલા ગણ તરફ પાછા જોઈને
તપાસી શકો છો કે તેને સંખ્યા $n$ \emph{તરીકે} જ ગણવામાં આવશે.

આ ઠરાવનું મહત્ત્વ \olref[sth][z][nat]{sec}માં ચર્ચીશું.
અત્યારે તે આપણને એક સહજ પરિણામ સાબિત કરવાની છૂટ આપે છે:

\begin{prop}\ollabel{naturalnumbersarentinfinite}
કોઈ પ્રાકૃતિક સંખ્યા ડેડેકિન્ડ-અનંત નથી.
\end{prop}

\begin{proof}
સાબિતી ગાણિતિક અનુમાનથી છે, એટલે કે
\olref[sfr][infinite][induction]{thm:dedinfiniteinduction}થી.
સ્પષ્ટ છે કે $0
= \emptyset$ ડેડેકિન્ડ-અનંત નથી. અનુમાનના પગલા માટે આપણે
પ્રતિવિધાન સ્થાપિત કરીશું: જો (અસંભવપણે) $s(n)$ ડેડેકિન્ડ-અનંત
હોય, તો $n$ ડેડેકિન્ડ-અનંત હોય.

તેથી ધારો કે $s(n)$ ડેડેકિન્ડ-અનંત છે, એટલે કે એવું !!{injection}
$f$ છે કે $\ran{f}\subsetneq \dom{f} = s(n) = n \cup
\{n\}$. બે કિસ્સા વિચારવાના છે.

\emph{કિસ્સો 1: $n \notin \ran{f}$.} તેથી $\ran{f} \subseteq n$,
અને $f(n) \in n$. $g = \funrestrictionto{f}{n}$ લો; હવે $\ran{g} =
\ran{f} \setminus \{f(n)\} \subsetneq n = \dom{g}$.
તેથી $n$ ડેડેકિન્ડ-અનંત છે.

\emph{કિસ્સો 2: $n \in \ran{f}$.} કોઈ નિશ્ચિત
$m \in \dom{f} \setminus \ran{f}$ લો, અને
$s(n) = n \cup \{n\}$ પ્રદેશ ધરાવતો વિધેય $h$ વ્યાખ્યાયિત કરો:
\[
h(x) =
	\begin{cases}
		f(x) & \text{જો }f(x) \neq n\\
		m & \text{જો }f(x)=n
	\end{cases}
\]
આમ $h(f^{-1}(n)) = m
\neq n = f(f^{-1}(n))$ સિવાય $h$ અને $f$ બધે એકસરખાં મૂલ્યો
આપે છે. $f$ એ !!a{injection} હોવાથી $n \notin
\ran{h}$; અને $\ran{h} \subsetneq \dom{h} = s(n)$.
હવે કિસ્સા 1ની દલીલથી $n$ ડેડેકિન્ડ-અનંત છે.
\end{proof}

જોકે અનંતતાની સ્વયંસિદ્ધિનું \emph{સમર્થન} કેવી રીતે કરવું તે
પ્રશ્ન બાકી રહે છે. ટૂંકો જવાબ એ છે કે આપણે જે વર્ણન આપતા આવ્યા
છીએ તેમાં બીજો સિદ્ધાંત ઉમેરવો પડશે. તે સિદ્ધાંત આ છે:
\begin{enumerate}
	\item[] \stagesinf. કોઈ અનંત તબક્કો છે. એટલે કે એવો તબક્કો છે
જે (ક) પહેલો તબક્કો નથી, અને જેની (ખ) પહેલાં અમુક તબક્કાઓ છે,
પણ જેને (ગ) તત્કાલ પૂર્વવર્તી તબક્કો નથી.
\end{enumerate}
આ સિદ્ધાંતથી અનંતતાની સ્વયંસિદ્ધિ સીધી મળે છે. આપણે જાણીએ છીએ
કે પ્રાકૃતિક સંખ્યા $n$ એ $n$મા તબક્કે રચાય છે. તેથી ગણ
$\omega$ પહેલા અનંત તબક્કે રચાય છે. અને $\omega$ પોતે
અનંતતાની સ્વયંસિદ્ધિનો સાક્ષી છે.

જોકે આનાથી આપણે ફરી \stagesinf{}નું સમર્થન કેવી રીતે કરવું તે
પ્રશ્ન પર પહોંચીએ છીએ. \stagessucc{}ની જેમ, આ સિદ્ધાંત પણ
સંચયી-પુનરાવર્તી પદાનુક્રમના આપેલા વર્ણનનો સ્પષ્ટ ભાગ નહોતો.
પણ વાત એટલી જ નથી: ગણો તબક્કાવાર રચાય એવા પુનરાવર્તી
પદાનુક્રમનો ખ્યાલ જ પ્રક્રિયામાં \emph{અનંત} તબક્કો સામેલ હોય
એવું માનવા આપણને ફરજ પાડતો નથી. તબક્કાઓનો ક્રમ પ્રાકૃતિક
સંખ્યાઓ જેવો હોય એમ માનવું પૂરેપૂરું સુસંગત લાગે છે.

જોકે આનાથી એક સ્પષ્ટ પ્રશ્ન ઊભો થાય છે.
\olref[sfr][infinite][dedekindsproof]{sec}માં આપણે
ડેડેકિન્ડ-અનંત (વિચારોનો) ગણ હોય એવી ડેડેકિન્ડની “સાબિતી”
વિચારી હતી. તમને તે બહુ સંતોષકારક ન લાગી હોય. પણ
\stagesinf{}ને પુનરાવર્તી ગણ-કલ્પના (અથવા “વિચારના નિયમો”)
આપણા પર “લાદતી” ન હોય, તો ડેડેકિન્ડ-અનંત ગણ હોય તે દાવાના
આંતરિક સમર્થન વિના આપણે હજી રહીએ છીએ.

અલબત્ત, અહીં ઘણું વધુ કહી શકાય. પણ હવે સ્વયંસિદ્ધિનું
(અથવા સિદ્ધાંતનું) સમર્થન કરવા \emph{શું જોઈએ} તે વિશે
વિચારવાનું શરૂ કરી શકો એવી સ્થિતિએ તમે પહોંચ્યા હશો.
આગળ આપણે \stagesinf{}ને સ્વીકૃત માનીશું.

\end{document}
